\documentclass[12pt,a4paper]{article}
\usepackage[utf8]{inputenc}
\usepackage[margin=1in]{geometry}
\usepackage{times}
\usepackage{setspace}
\onehalfspacing
\usepackage{titlesec}
\usepackage{enumitem}

\titleformat{\section}
{\Large\bfseries}
{}
{0em}
{}[\titlerule]

\titleformat{\subsection}
{\large\bfseries}
{}
{0em}
{}

\begin{document}

\begin{center}
{\Large\bfseries Documentation of Patents and Manuscript Briefs}\\[0.5cm]
{\large Deborah Akuoko}\\[0.3cm]
{\normalsize University of Edinburgh}\\[0.5cm]
\today
\end{center}

\vspace{0.5cm}

\section{Patents}

This section documents the patent applications resulting from my doctoral research at the University of Edinburgh.

\subsection{GB Patent Application 1: Spatiotemporal Fusion}

\textbf{Title:} Spatiotemporal Fusion Method and System

\textbf{Status:} Patent Application Filed

\textbf{Description:} This patent application covers innovative methods and systems for spatiotemporal data fusion, representing a significant contribution to the field of signal processing and data integration. The invention addresses the challenge of combining temporal and spatial information in a unified framework, with applications across multiple domains including computer vision, remote sensing, and data analytics.

\textbf{Key Contributions:}
\begin{itemize}[leftmargin=*]
\item Novel approach to spatiotemporal data fusion
\item Applications in computer vision and signal processing
\item Potential for commercialisation in various industries
\end{itemize}

\subsection{GB Patent Application 2: Food Safety}

\textbf{Title:} Food Safety Detection and Analysis System

\textbf{Status:} Patent Application Filed

\textbf{Description:} This patent application presents an advanced system for food safety detection and analysis, utilising innovative sensing and analysis techniques. The invention addresses critical challenges in food safety monitoring, providing enhanced capabilities for detection, analysis, and quality assurance in food production and distribution systems.

\textbf{Key Contributions:}
\begin{itemize}[leftmargin=*]
\item Advanced food safety detection methodology
\item Integration of sensing and analytical technologies
\item Applications in food production and quality assurance
\end{itemize}

\vspace{0.5cm}

\section{Manuscripts Briefs}

This section documents the working papers and manuscript briefs emerging from my doctoral research, demonstrating the scholarly contributions and intellectual property developed during my PhD studies.

\subsection{Paper 1: PAWD Framework}

\textbf{Title:} PAWD Framework: A Comprehensive Approach to Signal Processing

\textbf{Status:} Working Paper / Manuscript Brief

\textbf{Description:} This paper presents the PAWD (Processing, Analysis, Workflow, and Decision) framework, providing a comprehensive approach to signal processing and data analysis. The framework offers a structured methodology for handling complex data processing tasks across various domains.

\subsection{Paper 2: Material Detection}

\textbf{Title:} Advanced Material Detection Methods and Applications

\textbf{Status:} Working Paper / Manuscript Brief

\textbf{Description:} This paper explores innovative methods for material detection, presenting novel techniques and their applications. The research contributes to the field of sensing and detection technologies, with potential applications in quality control, security, and industrial processes.

\subsection{Paper 3: Spatiotemporal Fusion}

\textbf{Title:} Spatiotemporal Fusion: Methods and Applications

\textbf{Status:} Working Paper / Manuscript Brief

\textbf{Description:} This paper provides a detailed examination of spatiotemporal fusion techniques, building upon the patent application in this area. The work presents theoretical foundations, methodological approaches, and practical applications of spatiotemporal data fusion across multiple domains.

\subsection{Paper 4: Food Safety}

\textbf{Title:} Food Safety Detection and Analysis: A Comprehensive Approach

\textbf{Status:} Working Paper / Manuscript Brief

\textbf{Description:} This paper complements the food safety patent application, providing detailed analysis of detection methods, analytical techniques, and their applications in food safety monitoring. The research addresses critical challenges in food production and distribution systems.

\subsection{Paper 5: Complexity Analysis}

\textbf{Title:} Complexity Analysis in Signal Processing and Data Systems

\textbf{Status:} Working Paper / Manuscript Brief

\textbf{Description:} This paper examines complexity analysis methods in signal processing and data systems, providing insights into computational complexity, system performance, and optimisation strategies. The work contributes to understanding the theoretical and practical aspects of complex data processing systems.

\vspace{0.5cm}

\section{Summary}

The patents and manuscript briefs documented above represent substantial intellectual contributions emerging from my doctoral research at the University of Edinburgh. These works demonstrate:

\begin{itemize}[leftmargin=*]
\item \textbf{Innovation:} Novel approaches to spatiotemporal fusion and food safety detection
\item \textbf{Technical Expertise:} Advanced knowledge in signal processing, data analysis, and system design
\item \textbf{Research Impact:} Contributions to multiple fields including computer vision and food safety
\item \textbf{Commercial Potential:} Two patent applications indicating potential for commercialisation and practical application
\item \textbf{Scholarly Contribution:} Multiple working papers and manuscript briefs demonstrating research productivity and academic rigour
\end{itemize}

These achievements reflect my competence, expertise, and potential to make meaningful contributions to the academic and research landscape in the United Kingdom, as recognised by my supervisors and the University of Edinburgh.

\vspace{0.3cm}

\begin{flushright}
Deborah Akuoko\\
University of Edinburgh\\
\today
\end{flushright}

\end{document}

